\documentclass[a4paper,11pt]{article}

% Pacotes Principais -----------------------------------------------------------
\usepackage[portuges,brazil]{babel}
\usepackage[utf8]{inputenc}

% Figuras e Imagens ------------------------------------------------------------
\usepackage{graphicx}
% Figuras lado a lado
\usepackage{epsfig}
\usepackage{subfigure}

% Utilizar H para inserir as imagens REALMENTE onde eu desejo
\usepackage{float}

% Fontes -----------------------------------------------------------------------
\usepackage[T1]{fontenc}
\usepackage{pslatex}

% Simbolos ---------------------------------------------------------------------
\usepackage{textcomp}

% Tabelas ----------------------------------------------------------------------
%\usepackage{multicol}
\usepackage{multirow}
% Colorir a tabela
\usepackage{colortbl}

% Outros pacotes ---------------------------------------------------------------
\usepackage{noitemsep}

\usepackage{color}
\usepackage{xcolor}

% Comentários em bloco
\usepackage{verbatim}

% Sublinhado, traçado...
\usepackage{ulem}

% Desenhar
\usepackage{tikz}

% Comandos gerais --------------------------------------------------------------

% Configuração da fonte
\renewcommand{\familydefault}{\sfdefault}

% Comandos matemáticos ---------------------------------------------------------
% Implicação em fórmulas
\newcommand{\implica}{$\quad\Rightarrow\quad$} %Meio de linha
\newcommand{\implicafim}{$\quad\Rightarrow$}   %Fim de linha
\newcommand{\tende}{$\rightarrow$}

% Fração com parenteses
\newcommand{\pfrac}[2]{\parent{\frac{#1}{#2}}}

% Transformada de Laplace e transformada Z
\newcommand{\lapl}{\pounds}
\newcommand{\transfz}{\mathcal{Z}}

% Sequências
\newcommand{\sequencia}[4]{$#1_{#2}$, $#1_{#3}$, \ldots, $#1_{#4}$}

% Outros ----------------------------------------------------------------------
\newcommand{\chave}[1]{\left\{#1\right\}}
\newcommand{\colchete}[1]{\left[#1\right]}
\newcommand{\parent}[1]{\left(#1\right)}

\let\D\displaymath

\newtheorem{definicao}{Definição}
\newtheorem{exemplo}{Exemplo}
\newtheorem{lema}{Lema}
\newtheorem{observ}{Observação}
\newtheorem{teorema}{Teorema}

\newcommand{\defin}[1]{\begin{definicao}#1\end{definicao}}
\newcommand{\exemp}[1]{\begin{exemplo}#1\end{exemplo}}
\newcommand{\obs}[1]{\begin{observ}#1\end{observ}}

\newcommand{\pixel}{{\it pixel}}

% Tikz
\usetikzlibrary{matrix}


\pagestyle{empty}
\setlength{\parindent}{12mm}
%\usepackage{indentfirst}

\begin{document}

\begin{table}[H]
\centering
\begin{tabular}{cl}
% Primeira Linha
\multirow{7}{*}
{
\includegraphics[width=1.65cm]{imgs/uern}
} & \\
& Governo do Estado do Rio Grande do Norte\\
% Segunda Linha
& Secretaria de Estado da Educação, da Cultural e dos Desportos -- SECD\\
% Terceira Linha
& {\sc Universidade do Estado do Rio Grande do Norte -- UERN}\\
% Quarta Linha
& Pró-Reitoria de Ensino e Graduação -- PROEG\\
% Quinta Linha
& Ciências da Computação -- Trabalho prático para 2\textordfeminine\ unidade (10
pontos)\\
& 
\end{tabular}
\end{table}

\begin{center}
\Large
\uline{{\sc Trabalho Prático de Implementação}}
\end{center}

\section*{Informações gerais}

Vimos ao longo da disciplina que o processamento digital de informações é uma
área bastante ampla e que, consequentemente, abrange diversos assuntos. Sabendo
disso, propõe-se aqui um trabalho no qual os alunos deverão implementar um
sistema computacional que irá processar informações de áudio, vídeo ou imagem.
Tal trabalho será de grande importância para avaliar o conhecimento adquirido ao
longo do semestre e irá compor a nota da segunda unidade da disciplina.

O trabalho deverá ser entregue {\bf até às 23:59h da data estabelecida} em sala
de aula, podendo ser desenvolvido maneira individual ou em grupo de no máximo 4
(quatro) integrantes. Após estabelecida a data, \uline{{\bf não haverá
prorrogação do prazo}}.

A avaliação do trabalho será realizada em dois dias e todos os grupos deverão
estar presentes. A participação dos grupos, dando sugestões ou tirando dúvidas
sobre a implementação, também será levada em consideração na avaliação.

A nota será a mesma para todos os componentes do grupo, havendo penalização para
cada dia de atraso na entrega. Caso algum componente do grupo não possa
comparecer no dia da apresentação, a falta deverá ser justificada junto à
coordenação, de acordo com as normas da UERN. Nos casos em que a {\bf falta não}
for {\bf justificada}, o aluno receberá {\bf nota 0.0 (zero)}.

\section*{Os critérios de avaliação}
Os trabalhos poderão ser implementados em qualquer linguagem de programação
desde que esta possua recursos para o desenvolvimento de interfaces gráficas.
Sugere-se, entretanto, a utilização de C/C++ (Qt ou Gtk), Java ou Python. 

Os seguintes critérios serão levados em consideração para composição da nota
final:

\begin{enumerate}
    \item Apresentação do sistema (30\% da nota)
    \item Objetivos/Resultados (25\% da nota)
    \item Interface amigável e de fácil utilização (15\% da nota)
    \item Documentação/Organização do código (15\% da nota)
    \item Participação em outras apresentações (15\% da nota)
\end{enumerate}

Os arquivos referentes a implementação e a apresentação deverão ser organizados
em uma estrutura de pastas, compactados e enviados para {\tt diogolr@gmail.com}
com o Assunto {\tt [SistMult-Natal] Trabalho - Grupo X} para a turma do Campus
de Natal, ou {\tt [SistMult-NC] Trabalho - Grupo X} para a turma do Núcleo de
Nova Cruz.

\section*{Nota final}
A nota final será composta pela soma dos critérios acima estabelecidos e pela
penalização por dia de atraso, conforme Eq. \ref{eq:nota}.

\begin{equation}\label{eq:nota}
N_f = \D 1,25^{-D_a} \times \sum_{i = 1}^5 C_i
\end{equation}

\noindent em que $D_a$ se refere aos dias de atraso contados a partir da data
estabelecida e $C_i$ corresponde ao $i$-ésimo critério.

\section*{Tópicos para os trabalhos}

Dentre os diversos assuntos estudados, sugere-se os seguintes tópicos: 

\begin{itemize}
    \item Compressão de imagem
    \begin{itemize}
        \item Com perda
        \item Sem perda
    \end{itemize}
    \item Compressão de vídeo
    \item Compressão de áudio
    \item Manipulação de imagens
    \begin{itemize}
        \item Conversão de formato (não será permitido a utilização de
              bibliotecas)
        \item Reconhecimento de padrões
        \item Eliminação de ruído
        \item Aplicação de filtros
        \item Detecção de arestas
        \item Editor: Fazendo, no mínimo, recorte, colagem (composição), {\it
              clip} e salvando a imagem final (pode-se utilizar as bibliotecas
              disponíveis)
    \end{itemize}
    \item Manipulação de vídeo
    \begin{itemize}
        \item Conversão de formato (não será permitido a utilização de
              bibliotecas)
        \item Aplicação de efeitos
        \item Editor: Fazendo, no mínimo, recorte, {\it play} e salvando o vídeo
              final (pode-se utilizar as bibliotecas disponíveis)
    \end{itemize}
    \item Separação áudio/vídeo (pode-se utilizar bibliotecas)
    \item Biometria: Identificação biométrica (dedo, face ou mão)\footnote{A
          aquisição da imagem/características deve ser feita pelo sistema
          desenvolvido, podendo-se utilizar qualquer equipamento/{\it
          drivers} disponíveis.}
    \item Aplicações musicais: Utilizando arquivos MID (por exemplo), tal como
          um teclado virtual, ou um violão. A aplicação deverá permitir que
          sejam executadas notas sequenciais (música) a partir de uma entrada de
          texto
    \item Desenvolvimento de jogos, utilizando OpenGL
\end{itemize}

{\bf Os alunos não estão restritos à esses tópicos}. Qualquer outro tópico,
desde que relacionado com o processamento de imagens, áudio ou vídeo, poderá ser
enviado por e-mail até 20 dias antes da data da apresentação, estando sujeito à
aprovação pelo professor.

Após a escolha do tema pelo grupo, qualquer modificação só poderá ser realizada
até 20 dias antes da data da apresentação (através de e-mail).

\end{document}
